\documentclass[sigconf,nonacm]{acmart}

\usepackage{booktabs}
\usepackage{graphicx}
\usepackage{amsmath}
\usepackage{url}
\usepackage{adjustbox}

\settopmatter{printacmref=false}
\renewcommand\footnotetextcopyrightpermission[1]{}
\pagestyle{plain}

\setlength{\emergencystretch}{3em}
\begin{document}

\title{Single-Kernel Fusion for Sequential Fitness Evaluation via WebGPU Compute Shaders}

\author{Ahmet Baris Gunaydin}
\affiliation{%
  \institution{Independent Researcher}
  \city{Bolu}
  \country{Turkey}
}
\email{abgunaydin94@gmail.com}

\begin{abstract}
Fusing sequential fitness evaluations into single GPU compute shader dispatches eliminates the per-step kernel launch overhead that dominates framework-based GPU computation. We demonstrate this across \textbf{four GPU APIs on two hardware platforms}. On the same Tesla T4: a hand-fused CUDA kernel achieves 720$\times$ over PyTorch per-step dispatch on Acrobot-v1, JAX \texttt{lax.scan} achieves 172$\times$, and Triton achieves 27$\times$. On the same Apple M2 Pro: a WebGPU compute shader achieves 159$\times$ over PyTorch MPS on a 1,500-timestep financial simulation, and 54$\times$ on Acrobot. An unfused ablation isolates 2.18$\times$ from fusion alone. We show \texttt{torch.compile} fails at $L \geq 1{,}000$ and that the advantage scales with episode length $L$. A native Metal baseline via wgpu quantifies Chrome's browser overhead at 1.92$\times$. The insight --- hand-fused kernels dramatically outperform framework dispatch on sequential workloads --- is \textbf{GPU-API-agnostic}, and WebGPU makes such fusion zero-install.
\end{abstract}

\keywords{WebGPU, compute shaders, kernel fusion, neuroevolution, GPU computing, evolutionary computation, browser-based computing}

\maketitle

%% ═══════════════════════════════════════════════════════════
\section{Introduction}

GPU-accelerated evaluation of population-based algorithms requires dispatching the same computation across thousands of individuals per generation. When the fitness function is embarrassingly parallel (e.g., Rastrigin), modern ML frameworks like PyTorch and JAX efficiently batch these evaluations. However, many real-world fitness functions contain \textbf{sequential dependencies} --- reinforcement learning rollouts, financial simulations, control system evaluations --- where each timestep depends on the previous state.

For these sequential workloads, framework-based approaches face a fundamental bottleneck: \textbf{per-step kernel dispatch overhead}. PyTorch's eager execution model requires a CPU$\to$GPU round-trip for each of the $L$ timesteps within each generation, totaling ${\sim}15 \times L$ kernel launches per generation. For $L = 1{,}500$, this produces ${\sim}22{,}500$ GPU operations per generation with CPU$\leftrightarrow$GPU synchronization between each.

We present a solution: \textbf{single-kernel fusion}, where the entire multi-step evaluation (all $L$ timesteps for all $N$ individuals) executes as a single GPU compute shader dispatch. The kernel maintains per-individual state in GPU memory and iterates through timesteps entirely on-device, eliminating all intermediate host synchronization.

This approach is implemented in WGSL (WebGPU Shading Language~\cite{w3c2023}), providing: (1)~zero installation --- the compute shader runs in any WebGPU-capable browser, and (2)~cross-vendor portability --- the same shader executes on Apple Metal, NVIDIA Vulkan, AMD Vulkan, and Intel D3D12 backends.

\textbf{Contributions:}
\begin{itemize}
  \item \textbf{C1: Kernel fusion for sequential fitness functions.} On the same M2 Pro GPU, fusing a 1,500-timestep financial simulation yields 159$\times$ over PyTorch MPS. On the same Tesla T4, a hand-fused CUDA kernel achieves 720$\times$ over PyTorch CUDA on Acrobot, JAX \texttt{lax.scan} achieves 172$\times$, and Triton achieves 27$\times$ --- all confirming the fusion advantage is not WebGPU-specific. An unfused ablation isolates 2.18$\times$ from fusion alone.
  \item \textbf{C2: Zero-install GPU compute via WebGPU.} The complete engine (${\sim}350$ lines of WGSL) runs on any WebGPU-capable browser.
  \item \textbf{C3: Empirical characterization} of throughput scaling, numerical precision, browser sandbox overhead (native Metal baseline), NVIDIA CUDA comparison, and JAX GPU comparison.
\end{itemize}

%% ═══════════════════════════════════════════════════════════
\section{Related Work}

\textbf{GPU-based evolutionary computation.} Langdon~\cite{langdon2010} surveys GPU implementations of GP in CUDA. DEAP~\cite{fortin2012} provides a Python EA framework. EvoJAX~\cite{tang2022} achieves hardware-accelerated neuroevolution via JAX on TPU/GPU. EvoGP~\cite{wu2025} and GSGP-CUDA~\cite{trujillo2022} provide CUDA-based GP frameworks. Luong et al.~\cite{luong2010} demonstrated GPU-based island models. All require CUDA/TPU infrastructure.

\textbf{WebGPU for computation.} WGPY~\cite{hidaka2025} provides a NumPy-like array library via WebGPU. Sengupta et al.~\cite{sengupta2025} benchmark WebGL vs WebGPU for browser-based EC but benchmark individual operations rather than full evolutionary loops.

\textbf{Browser-based EC.} Duda \& Dlubacz~\cite{duda2013b} demonstrated browser EC with JavaScript on CPU; they also explored WebGL acceleration~\cite{duda2013a}. Merelo-Guervos and Garcia-Sanchez~\cite{merelo2015} modeled distributed browser EC systems.

%% ═══════════════════════════════════════════════════════════
\section{System Architecture}

\subsection{GPU-Native Evolution Engine}

The evolutionary loop runs within two WGSL compute shaders: \texttt{nn\_core.wgsl} (shared RNN forward pass, 246 parameters) and \texttt{swarm\_shader.wgsl} (evaluate and evolve entry points). Genome data resides entirely in GPU VRAM using a ping-pong double-buffer pattern. The only CPU$\to$GPU transfer per generation is a 16-byte uniform buffer.

\subsection{Sequential Workload: Financial Domain}

Each individual is evaluated over $L = 1{,}500$ timesteps of historical cryptocurrency price data (16 assets, hourly candles, ${\sim}62$ days). The simulation models leveraged portfolio management with continuous allocation, asymmetric fees, liquidity impact, and maintenance margin liquidation. Each timestep depends on the previous portfolio state.

\subsection{Experimental Methodology}

WebGPU throughput: automated Puppeteer benchmarks ($N = 10$ runs, 20s warmup). Python baselines: $N = 10$ runs, mean $\pm$ std. Hardware: Apple M2 Pro (Chrome 123, macOS 14) and Tesla T4 (CUDA 12.8, Google Colab).

%% ═══════════════════════════════════════════════════════════
\section{Results}

\subsection{Embarrassingly Parallel Workload (Rastrigin)}

\begin{table}[h]
\caption{Rastrigin~\cite{rastrigin1974} Benchmark (POP=4,096, DIM=2,000)}
\label{tab:rastrigin}
\adjustbox{max width=\columnwidth}{\small
\begin{tabular}{lrrr}
\toprule
System & gen/s & vs NumPy & $N$ \\
\midrule
Python/NumPy (CPU) & 3.9 $\pm$ 0.3 & --- & 10 \\
JAX jit (CPU) & 20.6 $\pm$ 0.5 & 5.3$\times$ & 10 \\
PyTorch (MPS, M2 Pro) & 160.5 $\pm$ 3.6 & 41$\times$ & 10 \\
\textbf{WebGPU (Chrome, M2 Pro)} & \textbf{170.3 $\pm$ 8.4} & \textbf{44$\times$} & 10 \\
wgpu-native (Metal, M2 Pro) & 326.5 $\pm$ 0.3 & 84$\times$ & 10 \\
PyTorch (CUDA, Tesla T4) & 311.1 $\pm$ 0.8 & 80$\times$ & 10 \\
JAX jit+vmap (CUDA, T4) & 1,163.9 $\pm$ 168.2 & 298$\times$ & 10 \\
\bottomrule
\end{tabular}}
\end{table}

\textbf{JAX GPU dominates} at 1,164~gen/s --- 6.8$\times$ faster than WebGPU. This is expected: XLA compilation on NVIDIA hardware represents the ceiling for parallel workloads. WebGPU achieves 1.06$\times$ parity with PyTorch MPS.

\textbf{Native Metal baseline.} The identical WGSL shader via wgpu-native (Rust) achieves 326.5~gen/s --- \textbf{1.92$\times$ faster than WebGPU in Chrome}, quantifying browser overhead at 48\%. PyTorch MPS (160.5) is slower than WebGPU (170.3) despite being native.

\subsection{Sequential Workload: Financial Simulation}

\begin{table}[h]
\caption{Financial Simulation (POP=10,000, 246 params, 1,500 timesteps)}
\label{tab:financial}
\adjustbox{max width=\columnwidth}{\small
\begin{tabular}{lrrrr}
\toprule
System & gen/s & vs NumPy & vs MPS & vs CUDA \\
\midrule
Python/NumPy (CPU) & 0.056 & --- & --- & --- \\
JAX jit+lax.scan (CPU) & 0.157 & 2.8$\times$ & --- & --- \\
PyTorch (MPS) & 0.29 & 5.2$\times$ & --- & --- \\
PyTorch (CUDA, T4) & 0.49 & 8.8$\times$ & 1.7$\times$ & --- \\
JAX lax.scan+vmap (T4) & 6.43 & 115$\times$ & 22$\times$ & --- \\
\textbf{WebGPU (Chrome)} & \textbf{46.2 $\pm$ 1.0} & \textbf{825$\times$} & \textbf{159$\times$} & \textbf{94$\times$} \\
\bottomrule
\end{tabular}}
\end{table}

\textbf{Same-hardware comparison (M2 Pro):} WebGPU achieves 159$\times$ over PyTorch MPS --- same GPU, same unified memory, isolating pure dispatch overhead. \textbf{Same-hardware comparison (T4):} JAX with \texttt{lax.scan}+\texttt{vmap} achieves 13$\times$ over PyTorch CUDA --- same PCIe path, isolating XLA fusion benefit. \textbf{Cross-hardware:} WebGPU on M2 Pro achieves 94$\times$ over PyTorch CUDA on T4 and 7.2$\times$ over JAX GPU on T4; these numbers include both fusion gains and M2 Pro's unified memory advantage, and should not be interpreted as pure algorithmic speedups. The cleanest hierarchy uses same-hardware pairs: hand-fused shaders (46.2, M2) $>$ XLA fusion (6.43, T4) $>$ per-step dispatch (0.29 MPS / 0.49 CUDA).

\subsection{Sequential Workload: Acrobot-v1}

\begin{table}[h]
\caption{Acrobot-v1 (POP=4,096, 163 params, 500 steps, RK4)}
\label{tab:acrobot}
\adjustbox{max width=\columnwidth}{\small
\begin{tabular}{llrrr}
\toprule
System & Hardware & gen/s & vs PyTorch & $N$ \\
\midrule
\multicolumn{5}{l}{\textit{Same hardware: Tesla T4 (CUDA)}} \\
PyTorch per-step & T4 & 0.61 $\pm$ 0.03 & 1$\times$ & 10 \\
Triton fused & T4 & 16.4 $\pm$ 1.0 & 27$\times$ & 10 \\
JAX lax.scan+vmap & T4 & 105.1 $\pm$ 1.0 & 172$\times$ & 10 \\
\textbf{Hand-fused CUDA kernel} & \textbf{T4} & \textbf{439 $\pm$ 121} & \textbf{720$\times$} & 10 \\
\midrule
\multicolumn{5}{l}{\textit{Same hardware: Apple M2 Pro (Metal)}} \\
PyTorch MPS per-step & M2 Pro & 2.52 $\pm$ 0.04 & 1$\times$ & 10 \\
wgpu-native fused & M2 Pro & 30.5 $\pm$ 0.4 & 12$\times$ & 10 \\
WebGPU unfused (Chrome) & M2 Pro & 62.3 $\pm$ 0.1 & 25$\times$ & 10 \\
\textbf{WebGPU fused (Chrome)} & \textbf{M2 Pro} & \textbf{135.9 $\pm$ 4.0} & \textbf{54$\times$} & 10 \\
\bottomrule
\end{tabular}}
\end{table}

\textbf{Generalizability confirmed.} On the same T4, four independent approaches all show massive fusion gains over PyTorch per-step dispatch: hand-fused CUDA (720$\times$), JAX XLA fusion (172$\times$), and Triton (27$\times$). On the same M2 Pro, WebGPU fused (54$\times$), WebGPU unfused (25$\times$), and wgpu-native (12$\times$) all outperform PyTorch MPS. \textbf{The fusion advantage is not WebGPU-specific --- it manifests across CUDA, Metal, XLA, and Triton.}

\textbf{Ablation (same hardware, M2 Pro).} Unfused WebGPU (500 separate dispatches) achieves 62.3~gen/s; fused achieves 135.9~gen/s --- a clean \textbf{2.18$\times$ from fusion alone}. On T4, hand-fused CUDA (439) vs Triton (16.4) shows that kernel quality matters enormously even among fused approaches.

\textbf{$L$-scaling.} JAX GPU achieves 105.1~gen/s at $L=500$ but only 6.43~gen/s at $L=1{,}500$ (financial), suggesting compiled fusion overhead scales with episode length while hand-fused kernels maintain constant overhead.

\subsection{torch.compile Scaling Failure}

\begin{table}[h]
\caption{\texttt{torch.compile} scaling (POP=10,000)}
\label{tab:compile}
\adjustbox{max width=\columnwidth}{\small
\begin{tabular}{lrrrr}
\toprule
Timesteps & Eager & Compiled & Compile time & Speedup \\
\midrule
10 & 0.067s & 0.002s & 1.8s & 30.7$\times$ \\
100 & 0.024s & 0.015s & 4.7s & 1.5$\times$ \\
500 & 0.158s & 0.081s & 25.5s & 2.0$\times$ \\
1,000 & 0.227s & \textbf{RecursionError} & --- & --- \\
5,000 & 3.4s & \textbf{OOM killed} & $>$30 min & --- \\
\bottomrule
\end{tabular}}
\end{table}

\subsection{Numerical Precision (f32 vs f64)}

Maximum absolute error in NN outputs: $1.97 \times 10^{-7}$. Fitness rankings perfectly preserved: Spearman $\rho = 1.000$ across 10,000 genomes.

\subsection{Scaling Characterization}

\begin{table}[h]
\caption{Population Scaling (Rastrigin, DIM=2,000)}
\label{tab:popscale}
\begin{tabular}{rrr}
\toprule
Population & gen/s & VRAM (MB) \\
\midrule
512 & 12,948 & 7.8 \\
1,024 & 12,791 & 15.6 \\
2,048 & 12,842 & 31.3 \\
4,096 & 12,685 & 62.5 \\
8,192 & 12,714 & 125.1 \\
16,384 & 12,741 & 250.1 \\
32,768 & 12,702 & 500.3 \\
\bottomrule
\end{tabular}
\end{table}

Throughput flat from 512 to 32,768 ($<$2.1\% variation). All five benchmark functions (Sphere, Rastrigin, Ackley, Schwefel, Griewank) within 8\%. Genome scaling (DIM 100--4,000) shows $<$2\% variation.

\subsection{GPU Utilization}

GPU utilization reaches 79.7\% $\pm$ 1.5\% during evolution with no thermal throttling (30.7\textdegree C constant, CV$<$5\% over 60s).

\subsection{Multi-Tab GPU Contention}

With 8 concurrent tabs, total throughput reaches 6.4$\times$ with 80\% per-tab efficiency.

\subsection{MountainCar-v0}

\begin{table}[h]
\caption{MountainCar-v0 (POP=4,096, 51 params, 200 timesteps)}
\label{tab:mountaincar}
\begin{tabular}{lrrr}
\toprule
System & gen/s & vs MPS & $N$ \\
\midrule
PyTorch (MPS, M2 Pro) & 18.7 $\pm$ 0.6 & --- & 10 \\
\textbf{WebGPU (Chrome)} & \textbf{1,258.8 $\pm$ 11.4} & \textbf{67$\times$} & 10 \\
\bottomrule
\end{tabular}
\end{table}

Third standard Gym environment. The 67$\times$ speedup at $L=200$ is consistent with the pattern.

\subsection{N-Body Simulation}

\begin{table}[h]
\caption{N-Body (512 bodies, 200 timesteps, $O(N^2)$ force, sequential)}
\label{tab:nbody}
\adjustbox{max width=\columnwidth}{\small
\begin{tabular}{lrrr}
\toprule
System & gen/s & vs MPS & $N$ \\
\midrule
PyTorch (CUDA, T4) & 15.9 $\pm$ 2.3 & --- & 10 \\
PyTorch (MPS, M2 Pro) & 30.5 $\pm$ 0.2 & --- & 10 \\
\textbf{WebGPU (Chrome, M2 Pro)} & \textbf{84.0 $\pm$ 0.2} & \textbf{2.8$\times$} & 10 \\
\bottomrule
\end{tabular}}
\end{table}

N-Body is compute-bound ($O(N^2)$ force per step) rather than dispatch-bound, yielding a smaller 2.8$\times$ speedup. This confirms the fusion advantage scales with dispatch overhead fraction, not raw compute.

\subsection{Monte Carlo Pi (Parallel Baseline)}

\begin{table}[h]
\caption{Monte Carlo Pi (4,096 workers $\times$ 100K samples, parallel)}
\label{tab:montecarlo}
\adjustbox{max width=\columnwidth}{\small
\begin{tabular}{lrrr}
\toprule
System & gen/s & vs MPS & $N$ \\
\midrule
PyTorch (MPS, M2 Pro) & 7.8 $\pm$ 0.4 & --- & 10 \\
PyTorch (CUDA, T4) & 11.5 $\pm$ 0.0 & --- & 10 \\
\textbf{WebGPU (Chrome, M2 Pro)} & \textbf{115.9 $\pm$ 0.6} & \textbf{14.9$\times$} & 10 \\
\bottomrule
\end{tabular}}
\end{table}

On this embarrassingly parallel workload, WebGPU achieves 14.9$\times$ over PyTorch MPS and 10.1$\times$ over PyTorch CUDA --- larger than the Rastrigin parallel gap (1.06$\times$). The difference is that Monte Carlo is memory-bound (409M random number generations), where WebGPU's shader-native RNG avoids PyTorch's \texttt{torch.rand} overhead.

\subsection{CartPole-v1 Validation}

CartPole-v1 solved in 76 $\pm$ 46 ms ($N = 30$, 97\% solve rate). \textbf{Caveat:} with 4,096 parallel rollouts, the system solves on generation 0 or 1, validating random search at scale rather than the evolutionary algorithm.

%% ═══════════════════════════════════════════════════════════
\section{Limitations and Future Work}

\begin{itemize}
  \item \textbf{Cross-hardware confounding.} WebGPU runs on M2 Pro (unified memory); PyTorch CUDA and JAX GPU run on T4 (discrete, PCIe). Cross-hardware ratios (94$\times$, 223$\times$, 7.2$\times$) include both algorithmic and architectural differences. Same-hardware ratios (159$\times$ on M2 Pro, 13$\times$ on T4, 2.18$\times$ fusion ablation) are the cleanest comparisons. Running WebGPU on T4 via headless Chrome was attempted but WebGPU is unavailable in headless environments.
  \item \textbf{NVIDIA tested on Tesla T4 only.} A100/H100 testing would characterize whether the advantage persists at the high end.
  \item \textbf{JAX GPU tested on T4 only.} On TPU or A100, JAX's XLA compiler may further close the gap.
  \item \textbf{Browser overhead quantified but not decomposed.} The 1.92$\times$ gap includes process isolation, IPC, and Dawn translation.
\end{itemize}

%% ═══════════════════════════════════════════════════════════
\section{Reproducibility}

Browser: Chrome 120+. OS: macOS 14+, Linux (Colab). GPU: Apple Metal (M2 Pro), NVIDIA CUDA (Tesla T4). Code: \url{https://github.com/abgnydn/webgpu-kernel-fusion}.

%% ═══════════════════════════════════════════════════════════
\section{Conclusion}

Hand-fused kernels outperform framework-based dispatch on sequential fitness functions across all tested GPU APIs. On the same Tesla T4, a hand-fused CUDA kernel achieves 720$\times$ over PyTorch per-step, JAX XLA fusion achieves 172$\times$, and Triton achieves 27$\times$. On the same M2 Pro, WebGPU achieves 54$\times$ over PyTorch MPS. The advantage is GPU-API-agnostic and scales with episode length $L$. On parallel workloads, JAX GPU dominates (6.8$\times$ over WebGPU), confirming the advantage is specific to sequential dependency chains.

WebGPU makes kernel fusion accessible with zero installation across Apple Metal, NVIDIA Vulkan, AMD Vulkan, and Intel DirectX backends. A native Metal baseline quantifies Chrome's browser overhead at 1.92$\times$.

%% ═══════════════════════════════════════════════════════════
\section*{Acknowledgments}

Drafting assistance was provided by Claude (Anthropic). All experimental results, benchmarks, code, and verification were performed by the author.

\bibliographystyle{ACM-Reference-Format}
\begin{thebibliography}{13}

\bibitem{langdon2010}
W.B. Langdon, ``Large scale bioinformatics data mining with parallel genetic programming on graphics processing units,'' in \emph{Parallel and Distributed Computational Intelligence}, SCI 269, Springer, 2010.

\bibitem{fortin2012}
F.A. Fortin, F.M. De Rainville, M.A. Gardner, M. Parizeau, and C. Gagn\'{e}, ``DEAP: Evolutionary algorithms made easy,'' \emph{JMLR}, vol.~13, 2012.

\bibitem{w3c2023}
W3C, ``WebGPU Specification,'' 2023. [Online]. Available: \url{https://www.w3.org/TR/webgpu/}

\bibitem{tang2022}
Y. Tang, Y. Tian, and D. Ha, ``EvoJAX: Hardware-accelerated neuroevolution,'' in \emph{Proc. GECCO Companion}, 2022.

\bibitem{rastrigin1974}
L.A. Rastrigin, \emph{Systems of Extremal Control}. Moscow: Nauka, 1974.

\bibitem{sengupta2025}
S. Sengupta, N. Wu, M. Varvello, K. Jana, and S. Chen, ``From WebGL to WebGPU: A reality check of browser-based GPU acceleration,'' in \emph{Proc. ACM Internet Measurement Conference}, 2025.

\bibitem{merelo2015}
J.J. Merelo-Guervos and P. Garcia-Sanchez, ``Modeling browser-based distributed evolutionary computation systems,'' \emph{arXiv:1503.06424}, 2015.

\bibitem{duda2013a}
J. Duda and W. Dlubacz, ``GPU acceleration for the web browser based evolutionary computing system,'' in \emph{Proc. ICSTCC}, 2013.

\bibitem{duda2013b}
J. Duda and W. Dlubacz, ``Distributed evolutionary computing system based on web browsers with JavaScript,'' in \emph{PARA 2012}, LNCS 7782, 2013.

\bibitem{luong2010}
T.V. Luong, N. Melab, and E.-G. Talbi, ``GPU-based island model for evolutionary algorithms,'' in \emph{Proc. GECCO}, 2010.

\bibitem{wu2025}
Z. Wu, L. Wang, K. Sun, Z. Li, and R. Cheng, ``Enabling population-level parallelism in tree-based genetic programming for GPU acceleration,'' \emph{arXiv:2501.17168}, 2025.

\bibitem{trujillo2022}
L. Trujillo, J.M. Munoz Contreras, D.E. Hernandez, M. Castelli, and J.J. Tapia, ``GSGP-CUDA: A CUDA framework for geometric semantic genetic programming,'' \emph{SoftwareX}, vol.~18, 2022.

\bibitem{hidaka2025}
M. Hidaka and T. Harada, ``WgPy: GPU-accelerated NumPy-like array library for web browsers,'' \emph{arXiv:2503.00279}, 2025.

\end{thebibliography}

\end{document}
